\documentclass[12pt,a4paper]{article}
\usepackage[utf8]{inputenc}
\usepackage[T2A]{fontenc}
\usepackage[russian]{babel}
\usepackage{geometry}
\usepackage{xcolor}
\usepackage{tcolorbox}
\usepackage{amsmath}
\usepackage{amssymb}

\geometry{margin=2cm}

\definecolor{smartgreen}{RGB}{0,100,0}
\definecolor{simpleblue}{RGB}{0,0,150}

\newcommand{\term}[3]{
\noindent\rule{\textwidth}{1pt}
\subsection*{#1}

\textcolor{smartgreen}{\textbf{УМНАЯ ВЕРСИЯ:}}

#2

\vspace{0.5em}
\textcolor{simpleblue}{\textbf{ДЛЯ ЛЮДЕЙ:}}

#3

\vspace{1em}
}

\begin{document}

\begin{center}
{\LARGE\textbf{СЛОВАРЬ ТЕРМИНОВ}}\\[0.5em]
{\large Умная версия + Объяснение для людей}
\end{center}

\vspace{1em}

Этот файл содержит все сложные термины из диплома с двумя объяснениями:
\begin{enumerate}
\item Умная версия — как написано в дипломе
\item Объяснение для людей — простыми словами
\end{enumerate}

\vspace{1em}

%% ========== ТЕРМИН 1 ==========
\term{1. РЕИДЕНТИФИКАЦИЯ (Re-ID)}
{Задача реидентификации формулируется как поиск по изображению-запросу $q$ из одной камеры соответствующих изображений того же физического экземпляра в галерее $\mathcal{G}$ из других камер. Необходимо построить отображение $f_\theta: \mathcal{X}\rightarrow\mathbb{R}^d$, переводящее изображение в эмбеддинг так, чтобы расстояния $D(f_\theta(q), f_\theta(g))$ были малы для $id(g)=id(q)$ и велики для иных $g$.}
{Представь: ты стоишь у входа в торговый центр и видишь красную машину. Потом идёшь на парковку и хочешь найти ту же самую машину среди сотни других. Вот это и есть реидентификация — найти одну и ту же тачку на разных камерах.

Почему это сложно? Потому что на одной камере машина может быть снята спереди, а на другой — сзади. Освещение разное, качество разное. А номер мы специально не используем (типа его не видно или заляпан грязью).}

%% ========== ТЕРМИН 2 ==========
\term{2. ЭМБЕДДИНГ (Embedding)}
{Отображение $f_\theta: \mathcal{X} \rightarrow \mathbb{R}^d$, переводящее изображение в вектор признаков в $d$-мерном евклидовом пространстве. В работе $d=2048$.}
{Эмбеддинг — это когда картинку с машиной превращаешь в набор чисел (вектор).

Представь, что каждую машину можно описать числами: <<краснота = 0.8, размер = 0.6, округлость = 0.3...>> и так далее 2048 чисел. Получается такой <<паспорт>> машины в виде чисел.

Зачем? Потому что сравнивать числа компьютеру проще, чем картинки. Если у двух картинок похожие числа — значит это одна и та же машина.}

%% ========== ТЕРМИН 3 ==========
\term{3. КРОСС-ЭНТРОПИЯ С ОСЛАБЛЕНИЕМ МЕТОК (Cross-Entropy with Label Smoothing)}
{Функция потерь для классификации с регуляризацией:
\[
\mathcal{L}_{CE} = -\sum_{i} q_i \log(p_i), \quad \text{где } q_i = (1-\varepsilon) \cdot y_i + \frac{\varepsilon}{K}
\]
При $\varepsilon=0.1$ истинному классу присваивается вероятность 0.9, остальным — $0.1/(K-1)$.}
{Окей, это про обучение нейросети.

\textbf{Обычная кросс-энтропия:} ты показываешь сети фотку и говоришь <<это машина №42, запомни>>. Сеть учится отвечать <<100\% это машина №42>>.

\textbf{Проблема:} сеть становится слишком самоуверенной. Она на 100\% уверена даже когда ошибается.

\textbf{Решение (Label Smoothing):} вместо <<это 100\% машина №42>> говорим <<это 90\% машина №42, а ещё может быть чуть-чуть похожа на другие>>.

Зачем? Сеть становится скромнее и лучше работает на новых данных. Это как учить студента: <<Будь уверен в ответе, но не на 100\%, оставляй место для сомнений>>.}

%% ========== ТЕРМИН 4 ==========
\term{4. TRIPLET LOSS}
{Функция потерь для метрического обучения:
\[
\mathcal{L}_{triplet} = \max\left(0, m + D(f_\theta(a), f_\theta(p)) - D(f_\theta(a), f_\theta(n))\right)
\]
где $(a, p, n)$ — якорь, позитив, негатив; $m$ — margin (зазор), в работе $m=0.3$.}
{Это хитрый способ обучения через тройки картинок.

Берём три картинки:
\begin{itemize}
\item \textbf{Якорь (anchor)} — фотка машины №1
\item \textbf{Позитив (positive)} — другая фотка той же машины №1
\item \textbf{Негатив (negative)} — фотка какой-то другой машины №2
\end{itemize}

Говорим сети: <<Сделай так, чтобы якорь был БЛИЖЕ к позитиву, чем к негативу>>.

Формула проще: расстояние(якорь, позитив) + запас < расстояние(якорь, негатив)

Запас (margin = 0.3) — это типа <<не просто ближе, а ощутимо ближе>>.

Результат: сеть учится группировать фотки одной машины вместе, а разные машины — разводить подальше.}

%% ========== ТЕРМИН 5 ==========
\term{5. BNNeck (Batch Normalization Neck)}
{Компонент архитектуры, нормализующий эмбеддинги перед функциями потерь:
\[
\hat{z} = \frac{\mathrm{BN}(z)}{\lVert \mathrm{BN}(z) \rVert_2}
\]
Позволяет эффективно совмещать классификационный и метрический подходы.}
{BNNeck — это такая <<прослойка>> в нейросети перед финальным шагом.

\textbf{Проблема:} у нас две задачи одновременно — классификация (определить ID машины) и метрическое обучение (сделать похожие машины близкими в пространстве). Они немного конфликтуют.

\textbf{Решение:} ставим BNNeck — он нормализует числа (приводит к единому масштабу) и позволяет обеим задачам работать вместе, не мешая друг другу.

\textbf{Аналогия:} это как переводчик между двумя людьми, которые говорят на разных языках. BNNeck переводит данные в формат, который подходит обеим задачам.}

%% ========== ТЕРМИН 6 ==========
\term{6. PK-СЕМПЛИНГ (PK Sampling)}
{Стратегия формирования батча: выбираются $P$ идентичностей, для каждой — $K$ изображений.
Батч $= P \times K$ изображений. В работе: $P=16$, $K=4$, итого $B=64$.}
{Это про то, как выбирать картинки для обучения.

\textbf{Обычный способ:} берём случайные 64 картинки из базы.

\textbf{PK-семплинг:} берём 16 разных машин ($P=16$), и для каждой машины берём 4 фотки ($K=4$). Итого $16 \times 4 = 64$ картинки.

Зачем? Чтобы в каждой порции данных были картинки одной и той же машины с разных ракурсов. Так Triplet Loss может находить хорошие тройки прямо внутри батча.

\textbf{Аналогия:} вместо того чтобы показывать сети 64 случайных фотки из альбома, показываем 16 мини-альбомов по 4 фотки каждого человека.}

%% ========== ТЕРМИН 7 ==========
\term{7. RANDOM ERASING}
{Аугментация: случайный прямоугольник изображения заполняется шумом или средним значением с вероятностью $p=0.5$. Повышает устойчивость к окклюзиям.}
{Это трюк при обучении: берём картинку и закрашиваем случайный кусок.

Зачем? Чтобы сеть не привязывалась к одной детали машины. Если закрыть фару — сеть всё равно должна узнать машину по другим признакам.

\textbf{Реальный пример:} на парковке часть машины может быть закрыта столбом или человеком. Сеть, обученная с Random Erasing, справится с этим лучше.

Это как тренировать футболиста играть с завязанным глазом — потом в реальной игре, когда оба глаза открыты, ему будет легче.}

%% ========== ТЕРМИН 8 ==========
\term{8. RE-RANKING (K-Reciprocal Re-ranking)}
{Постобработка результатов поиска с учётом взаимности ближайших соседей.
Если $A$ в топ-$k$ соседей $B$ и $B$ в топ-$k$ соседей $A$ — они <<взаимные соседи>>. Улучшает mAP на 5--7\% без переобучения.}
{Re-ranking — это перепроверка результатов после основного поиска.

\textbf{Обычный поиск:} нашли 10 машин, похожих на запрос. Вот они, отсортированы по похожести.

\textbf{Re-ranking:} а давайте проверим — если машина А похожа на машину Б, то Б тоже должна быть похожа на А, верно? Если это не так — что-то подозрительно.

\textbf{Пример:} ты ищешь друзей в соцсети. Обычный поиск: <<вот люди с похожими интересами>>. Re-ranking: <<а ты есть в их рекомендациях? Если нет — может, совпадение случайное>>.

Результат: точность поиска улучшается на 5--7\% БЕЗ дополнительного обучения. Просто умный пост-процессинг.}

%% ========== ТЕРМИН 9 ==========
\term{9. mAP (Mean Average Precision)}
{Средняя точность по всем запросам:
\[
\text{mAP} = \frac{1}{Q} \sum_{q} AP(q), \quad \text{где } AP(q) = \frac{1}{|\mathcal{R}_q|} \sum_{n=1}^{N} P@n \cdot \mathbf{1}\{r_n \in \mathcal{R}_q\}
\]
$P@n$ — precision на позиции $n$.}
{mAP — главная метрика качества поиска.

Объясню на примере. Ты ищешь свою машину. В базе 100 фоток, из них 5 — твоя машина.

Система выдала топ-10 результатов. Допустим, твоя машина на местах 1, 3, 7 (3 из 5 нашли в топ-10).

AP считает не просто <<сколько нашли>>, а <<насколько высоко в списке>>. Если все 5 твоих фоток в топ-5 — AP будет 1.0 (идеально). Если они разбросаны по всему списку — AP будет ниже.

mAP — это среднее AP по всем запросам. Чем выше — тем лучше система находит нужное.

У нас \textbf{mAP = 82.5\%} — это очень хорошо для этой задачи.}

%% ========== ТЕРМИН 10 ==========
\term{10. RANK-k (CMC — Cumulative Matching Characteristic)}
{Кумулятивная кривая совпадений. Rank-$k$ показывает вероятность того, что правильный ответ находится в топ-$k$ результатов поиска.}
{\textbf{Rank-1} = <<Система нашла правильную машину с первой попытки?>>

\textbf{Rank-5} = <<Правильная машина в первых 5 результатах?>>

\textbf{Rank-10} = <<Правильная машина в первых 10 результатах?>>

У нас \textbf{Rank-1 = 95.8\%} — это значит в 95.8\% случаев система с первого раза показывает правильную машину. Круто!

\textbf{Аналогия:} это как в поисковике Google. Rank-1 — ты нашёл нужное на первой строчке. Rank-10 — нужное где-то на первой странице.}

%% ========== ТЕРМИН 11 ==========
\term{11. COSINE DISTANCE / EUCLIDEAN DISTANCE}
{Метрики расстояния между векторами:
\begin{itemize}
\item Евклидово: $D(a,b) = \sqrt{\sum_i (a_i - b_i)^2}$
\item Косинусное: $D(a,b) = 1 - \frac{a \cdot b}{\|a\| \cdot \|b\|}$
\end{itemize}}
{Это способы измерить <<похожесть>> двух машин по их эмбеддингам.

\textbf{Евклидово расстояние:} обычное расстояние как линейкой. Если эмбеддинги — точки в пространстве, то это расстояние между ними.

\textbf{Косинусное расстояние:} измеряет угол между векторами. Не важно, насколько длинные вектора, важно куда они <<смотрят>>.

На практике: для нашей задачи оба работают, но косинусное чаще используется, потому что мы нормализуем эмбеддинги (делаем их одинаковой длины).

\textbf{Простая аналогия:} два человека показывают пальцем на объекты. Евклидово расстояние — насколько далеко их пальцы друг от друга. Косинусное — насколько они показывают в одном направлении.}

%% ========== ТЕРМИН 12 ==========
\term{12. WARMUP}
{Стратегия обучения: learning rate постепенно увеличивается от малого значения до целевого в течение первых $N$ эпох. В работе: 5 эпох warmup.}
{Warmup — это <<разогрев>> нейросети в начале обучения.

\textbf{Проблема:} если сразу начать учить с большой скоростью обучения (learning rate), сеть может <<сломаться>> — веса уйдут куда-то не туда.

\textbf{Решение:} первые 5 эпох learning rate постепенно растёт от почти нуля до нормального значения. Как спортсмен разминается перед тренировкой.

Зачем? Сеть успевает <<осмотреться>> в данных прежде чем начать активно учиться. Результат — более стабильное обучение.}

%% ========== ТЕРМИН 13 ==========
\term{13. COSINE ANNEALING SCHEDULER}
{Стратегия изменения learning rate по косинусоиде:
\[
lr(t) = lr_{min} + \frac{1}{2}(lr_{max} - lr_{min})\left(1 + \cos\left(\frac{\pi t}{T}\right)\right)
\]}
{Это про то, как менять скорость обучения по ходу тренировки.

\textbf{В начале:} учимся быстро (большой learning rate)

\textbf{К концу:} учимся медленно (маленький learning rate)

<<Cosine>> потому что график похож на косинус — плавно снижается от максимума к минимуму.

Зачем? В начале нужно быстро понять общую картину. В конце — точно настроить детали. Если учиться быстро до конца — промахнёшься мимо оптимума.

\textbf{Аналогия:} едешь на машине к цели. Сначала гонишь по трассе, потом замедляешься чтобы точно припарковаться.}

%% ========== ТЕРМИН 14 ==========
\term{14. ResNet-50}
{Остаточная сеть (Residual Network) с 50 слоями. Использует skip connections (остаточные связи) для решения проблемы затухающего градиента в глубоких сетях.}
{ResNet-50 — это очень популярная нейросеть для работы с картинками.

\textbf{<<50>>} = 50 слоёв. Чем больше слоёв — тем <<умнее>> сеть, но и сложнее обучать.

\textbf{<<Res>>} = residual (остаточный). Фишка: между слоями есть <<короткие пути>> (skip connections). Это как лифт в небоскрёбе — можно не идти по всем лестницам.

Зачем skip connections? Без них градиенты (сигналы обучения) затухают, пока идут через 50 слоёв. С ними — нормально доходят.

ResNet-50 — проверенный годами выбор. Не самый новый, но надёжный и быстрый.}

%% ========== ТЕРМИН 15 ==========
\term{15. HARD MINING}
{Стратегия выбора сложных примеров для обучения. Hard positive — далёкий позитив (большое $D(a,p)$), hard negative — близкий негатив (малое $D(a,n)$).}
{Hard mining — это выбор <<сложных>> примеров для обучения.

\textbf{Обычный подход:} берём случайные тройки (якорь, позитив, негатив).

\textbf{Hard mining:} специально ищем сложные случаи:
\begin{itemize}
\item \textbf{Hard positive:} фотка той же машины, но сеть думает что это другая (далеко в пространстве)
\item \textbf{Hard negative:} фотка другой машины, но сеть думает что это та же (близко в пространстве)
\end{itemize}

Зачем? Учиться на сложных примерах эффективнее. Нет смысла тренироваться на очевидных случаях — сеть и так их знает.

\textbf{Аналогия:} как готовиться к экзамену. Hard mining = решать задачи, которые не получаются, а не те, что уже умеешь.}

%% ========== ТЕРМИН 16 ==========
\term{16. DOMAIN SHIFT / DOMAIN ADAPTATION}
{Различие в распределении данных между камерами/датасетами. Domain Adaptation — методы адаптации модели к новому домену без полного переобучения.}
{Domain shift — это когда данные с разных камер выглядят по-разному.

\textbf{Проблема:} обучили сеть на камере А. Поставили на камеру Б — работает хуже. Потому что камера Б снимает с другими цветами, освещением, углом.

\textbf{Domain Adaptation} — методы, чтобы сеть работала на новых камерах без переобучения.

\textbf{Пример из жизни:} ты научился готовить на своей кухне. Приехал в гости — другая плита, другие кастрюли. Надо адаптироваться.

В нашей работе используем аугментации (Color Jitter) — это простой способ борьбы с domain shift.}

%% ========== ТЕРМИН 17 ==========
\term{17. YOLOv11 (You Only Look Once)}
{Однопроходной детектор объектов архитектуры YOLO версии 11 (Ultralytics, 2024). Предсказывает bbox и классы всех объектов за один прямой проход сети. YOLOv11n ($\approx$2.6M параметров) обеспечивает $\approx$100~fps на GPU при mAP\textsubscript{50}=39.5\% на COCO.}
{YOLO расшифровывается как «You Only Look Once» — смотришь только один раз. Обычные детекторы смотрят дважды: сначала ищут «интересные места», потом классифицируют. YOLO делает всё за один проход.

\textbf{Зачем нужен в нашей задаче?} Датасет VeRi-776 содержит уже готовые вырезанные машины. Но в реальной жизни камера видит весь перекрёсток целиком. YOLOv11 находит каждую машину на кадре и вырезает её. Потом этот кроп уже идёт в ReID-модель.

\textbf{Пайплайн:} камера → YOLOv11 (найти машину) → кроп → ViT-Small (узнать машину).

\textbf{Аналогия:} YOLOv11 — как охранник у входа, который замечает всех людей в толпе. ViT-Small — как детектив, который узнаёт конкретного человека в лицо.

\textbf{Результат в работе:} при использовании YOLO-кропа вместо GT-кропа Rank-1 падает с 90.05\% до 85.88\% — это <<цена>> реализма.}

%% ========== ТЕРМИН 18 ==========
\term{18. РАССТОЯНИЕ МАХАЛАНОБИСА}
{Метрика расстояния, учитывающая корреляционную структуру пространства:
\[
d_M(x, y) = \sqrt{(x-y)^\top \Sigma^{-1} (x-y)},
\]
где $\Sigma$ — ковариационная матрица, оцениваемая по обучающим данным. При $\Sigma = I$ совпадает с евклидовым расстоянием.}
{Обычное евклидово расстояние предполагает, что все оси одинаково важны. Но в реальных данных некоторые признаки сильно коррелируют (например, два признака цвета почти одинаковые).

\textbf{Аналогия:} представь карту города. Евклидово расстояние — прямая линия через здания. Махаланобис учитывает, что ехать по улице гораздо проще, чем через стены.

\textbf{Что случилось в нашей работе:} мы попробовали Махаланобис вместо L2 для поиска машин.
\begin{itemize}
    \item \textbf{ResNet-50}: почти ноль эффекта (+0.28\% mAP) — мы использовали диагональное приближение, что близко к обычной нормировке.
    \item \textbf{ViT-Small}: \emph{провал} — mAP упал с 69.4\% до 38.4\% (–31\%!).
\end{itemize}

\textbf{Почему провал?} ViT-эмбеддинги нормализуются на L2-норму и лежат на поверхности многомерной сферы. Косинусное расстояние — это угол между точками на сфере, что геометрически корректно. Махаланобис же «растягивает» пространство эллипсоидом, нарушая сферическую геометрию.

\textbf{Вывод:} для нормированных эмбеддингов косинусное расстояние оптимально. Это важный практический результат — попытка улучшить формулу метрики «сломала» систему.}

%% ========== ТЕРМИН 19 ==========
\term{19. ATTENTION MAP (Карта внимания)}
{Двумерная матрица весов $\mathbf{m} \in \mathbb{R}^{14\times14}$, показывающая, на какие пространственные патчи изображения обращает внимание CLS-токен ViT-Small. Агрегируется из последних $K=4$ трансформерных блоков операцией максимума по головам.}
{В ViT изображение делится на 196 патчей ($14\times14$). CLS-токен «смотрит» на все патчи через механизм внимания. Attention Map показывает, куда именно смотрит модель.

\textbf{Как выглядит:} тепловая карта поверх фото машины. Горячие (красные) зоны — туда модель смотрит больше. Холодные (синие) — игнорирует.

\textbf{Что мы обнаружили:}
\begin{itemize}
    \item При \emph{правильном} совпадении: модель концентрируется на уникальных деталях — номерной знак, рисунок диска, спойлер, логотип.
    \item При \emph{ошибке}: внимание размазано по форме и цвету — модель перепутала похожие машины.
\end{itemize}

\textbf{Важно:} модель \emph{не обучалась} определять «это фара» или «это диск» — у неё не было такой разметки. Но она сама научилась смотреть на дискриминативные детали!

\textbf{Аналогия:} как опытный полицейский сам знает, где у человека могут быть татуировки или шрамы — без специального обучения, просто из опыта.}

%% ========== ТЕРМИН 20 ==========
\term{20. PCA (Principal Component Analysis)}
{Метод снижения размерности: линейное преобразование данных в систему координат, где первые оси (главные компоненты) объясняют максимальную дисперсию. Для D-мерных эмбеддингов (D=384 или 2048) строится проекция на 2D для визуализации.}
{PCA — способ «сплющить» многомерные данные до 2D, чтобы можно было нарисовать.

\textbf{Задача:} каждый эмбеддинг машины — это 384 числа (у ViT). Нарисовать точку в 384-мерном пространстве невозможно. PCA находит лучший «угол обзора», с которого данные выглядят наиболее информативно, и проецирует на 2D.

\textbf{Что показали PCA-визуализации в дипломе:}
\begin{itemize}
    \item ResNet-50: кластеры одной машины размыты, перекрываются с другими.
    \item ViT-Small: кластеры компактные, хорошо разделены.
\end{itemize}

\textbf{Количественно}: коэффициент силуэта у ViT-Small = 0.26 против 0.024 у ResNet-50 — улучшение в 10 раз!

\textbf{Аналогия:} сфотографируй здание с лучшей стороны — с одного ракурса видно больше деталей. PCA выбирает этот лучший ракурс автоматически.}

%% ========== ТЕРМИН 21 ==========
\term{21. КОЭФФИЦИЕНТ СИЛУЭТА (Silhouette Score)}
{Метрика качества кластеризации: $s(i) = \frac{b(i) - a(i)}{\max(a(i), b(i))}$, где $a(i)$ — среднее расстояние до точек своего кластера, $b(i)$ — до ближайшего чужого кластера. Принимает значения от $-1$ до $+1$.}
{Это число, которое показывает насколько хорошо точки «знают, где их место».

\textbf{Смысл:} для каждой машины считается:
\begin{itemize}
    \item Насколько она близка к другим фоткам \emph{той же} машины?
    \item Насколько она далека от фоток \emph{других} машин?
\end{itemize}

Если ответ «близка к своим, далека от чужих» — силуэт близок к +1.\\
Если всё перемешалось — близок к 0.\\
Если попала в чужой кластер — отрицательный.

\textbf{Наши результаты:}
\begin{itemize}
    \item ResNet-50: силуэт = 0.024 (почти 0 — кластеры почти не разделены)
    \item ViT-Small: силуэт = 0.260 (в 10 раз лучше — кластеры чёткие)
\end{itemize}

Это математически подтверждает то, что мы видим визуально на PCA-картинках.}

%% ========== ТЕРМИН 22 ==========
\term{22. ГРАФ ПЕРЕХОДОВ МЕЖДУ КАМЕРАМИ}
{Направленный граф $G = (V, E)$, где $V$ — множество камер (20 штук), ребро $C_i \to C_j$ означает, что транспортное средство наблюдалось сначала камерой $i$, потом $j$. Вес ребра — число таких переходов.}
{Мы разобрали все маршруты 575 машин в обучающей выборке VeRi-776 и нарисовали граф: какие камеры смотрят на какие участки дороги.

\textbf{Что нашли:}
\begin{itemize}
    \item Граф содержит 108 из 380 возможных рёбер — дороги имеют чёткую топологию, не все камеры связаны.
    \item Камеры C11 и C06 — самые нагруженные (990 переходов). Скорее всего, стоят на одном перекрёстке.
    \item Средний маршрут машины: 8.9 камер (от 2 до 18).
    \item \emph{Проблемная камера C18}: Rank-1 = 0\% для всех пар с ней — камера снимает в совершенно других условиях.
\end{itemize}

\textbf{Практический вывод:} без анализа по камерам мы никогда не нашли бы камеру C18, которая «ломает» систему. В реальном видеонаблюдении это сигнал к техническому обслуживанию камеры.}

%% ========== ТЕРМИН 23 ==========
\term{23. FAISS (Facebook AI Similarity Search)}
{Библиотека для эффективного приближённого поиска ближайших соседей в многомерном пространстве. Поддерживает индексы IVF (Inverted File), PQ (Product Quantization), HNSW. Позволяет искать среди миллионов векторов за миллисекунды.}
{Когда машин в базе много — наивный перебор всех эмбеддингов медленный.

\textbf{Аналогия:} ты ищешь нужную книгу. Наивный подход — перебирать каждую из миллиона книг. FAISS — как каталог в библиотеке: сначала смотришь в алфавитный указатель, потом идёшь в нужный раздел.

\textbf{Как работает:}
\begin{enumerate}
    \item Векторы кластеризуются (IVF — создаётся <<оглавление>>).
    \item При поиске проверяется только ближайший кластер (а не все).
    \item Небольшая потеря точности в обмен на скорость в 1000+ раз.
\end{enumerate}

\textbf{В нашей задаче:} для real-time поиска на перекрёстке используется FAISS с косинусным расстоянием. Re-ranking (медленный, $O(N^2)$) — только для offline-криминалистики когда надо максимально точно найти машину из архива.}

%% ========== ТЕРМИН 24 ==========
\term{24. t-SNE (t-distributed Stochastic Neighbor Embedding)}
{Нелинейный метод снижения размерности для визуализации. В отличие от PCA, сохраняет локальную структуру данных: близкие точки в высокомерном пространстве остаются близкими на плоскости. Параметр perplexity контролирует число эффективных соседей.}
{t-SNE — более умный способ нарисовать многомерные данные по сравнению с PCA.

\textbf{PCA vs t-SNE:}
\begin{itemize}
    \item PCA: линейное сжатие, сохраняет глобальную структуру.
    \item t-SNE: нелинейное, сохраняет соседство — «кто рядом с кем».
\end{itemize}

\textbf{Зачем:} если в 384-мерном пространстве эмбеддинги одной машины образуют компактное облако — t-SNE покажет это облако как кучку точек на плоскости.

\textbf{В нашем случае:} вместо просто цветных точек мы рисуем на карте миниатюры изображений машин. Сразу видно, что модель считает похожим — похожие машины стоят рядом.

\textbf{Ограничение:} t-SNE не сохраняет масштаб. Расстояния между кластерами \emph{не означают} реальные расстояния в 384D пространстве.}

\end{document}
